%% -*- coding:utf-8 -*-
\author{Stefan Müller}
\title{Syntax der germanischen Sprachen}
\subtitle{Ein beschränkungsbasierter Ansatz\\\bigskip\rot{KI-übersetzt aus dem Englischen, Benutzung auf eigene Gefahr}}

\typesetter{Stefan Müller}


\BackTitle{Syntax der germanischen Sprachen}
\BackBody{Dieses Buch ist eine Einführung in die syntaktischen Strukturen der germanischen Sprachen. Die Analysen sind im Rahmen von HPSG light gehalten, einer vereinfachten Version von HPSG, die Bäume zur Darstellung der Analysen verwendet statt komplizierter Attribut-Wert-Matrizen. Das Buch richtet sich an Studierende mit Grundkenntnissen in den Bereichen Kasus, Konstituententests und einfache Phrasenstrukturgrammatiken (fortgeschrittenes BA- oder MA-Niveau) sowie an Forschende mit Interesse an den germanischen Sprachen und/oder an Head-Driven Phrase Structure Grammar/Sign-Based Construction Grammar, die jedoch nicht die Zeit haben, sich mit allen Details dieser Theorien zu beschäftigen.
}

\dedication{Für alle, die unsere Gesellschaften und Ökosysteme vor dem Kollaps zu bewahren versuchen}
%\renewcommand{\lsISBNdigital}{978-3-96110-408-6}
%\renewcommand{\lsISBNhardcover}{978-3-98554-066-2} Lehrbuch -> no hardcover
%\renewcommand{\lsISBNsoftcover}{978-3-98554-066-2}
%\renewcommand{\lsBookDOI}{10.5281/zenodo.7733033}
\renewcommand{\lsSeries}{tbls} % use lowercase acronym, e.g. sidl, eotms, tgdi
%\renewcommand{\lsSeriesNumber}{12} %will be assigned when the book enters the proofreading stage
%% \renewcommand{\lsURL}{http://langsci-press.org/catalog/book/25} % contact the coordinator for the right number
%\renewcommand{\lsID}{353}

\proofreader{%
}